% Part: incompleteness
% Chapter: representability-in-q
% Section: beta-function

\documentclass[../../../include/open-logic-section]{subfiles}

\begin{document}

\olfileid{inc}{req}{bet}
\olsection{بیٹا فنکشن دا معاون قضیہ}

سانوں وکھانا اے کہ جمع، ضرب تے $\Char{=}$ موجود ہون تاں
اسیں بنیادی بازگشت دا کم کر سکدے آں۔ ایس لئی سانوں سلسلیاں
اُتے کم کرن والے فنکشن بنانے پین گے۔ (جے قوت چڑھاؤن دا عمل وی
ہووے، تاں کم سوکھا ہووے۔) جدوں اسیں بنیادی بازگشت ورت سکدے ساں،
تاں «$n$-واں اولی عدد» ورگیاں شےواں دی تعریف کر سکدے ساں تے
کوڈ دین دا سِدھا طریقہ چُن سکدے ساں۔ ایتھے بنیادی بازگشت
پہلے توں ساڈے کول نہیں۔ اسیں وکھانا چاہندے آں کہ کمترین قدر
دی تلاش نال بنیادی بازگشت دا کم ہو سکدا اے۔
ایس لئی سانوں سلسلیاں دے کوڈاں لئی ہور طریقہ لبھنا پینا اے۔

\begin{lem}
\ollabel{lem:beta}
اک فنکشن $\beta(d,i)$ موجود اے جیہدے لئی ایہ شرط پوری ہوندی
اے۔ ہر سلسلے $a_0$، \dots،~$a_n$ لئی کوئی عدد~$d$ ایسا ملدا اے
کہ ہر $i \le n$ لئی $\beta(d,i) = a_i$ ہووے۔ نالے $\beta$ دی
تعریف بنیادی فنکشناں توں صرف ترکیب تے باقاعدہ کمترین قدر دی
تلاش نال کیتی جا سکدی اے۔
\end{lem}

$d$ نوں سلسلے $\tuple{a_0, \dots, a_n}$ دا کوڈ سمجھو۔ فنکشن
$\beta(d,i)$ ایس سلسلے دا $i$-واں رکن دیندا اے۔ (ایہ اوہ کوڈکاری
\emph{نہیں} جیہڑی اسیں اولی عدداں دیاں قوتاں نال پہلے کر چکے آں۔)
معاون قضیہ صرف تھوڑی جیہی گل ثابت کردا اے۔ ایہ سلسلے جوڑن یا
آخر وچ رکن پاؤن دا طریقہ نہیں دیندا۔ ایہ وی نہیں آکھدا کہ ترکیب
تے باقاعدہ کمترین قدر دی تلاش نال بنے فنکشناں نوں ورت کے~$d$
دا $a_0$، \dots،~$a_n$ توں \emph{حساب} لایا جا سکدا اے۔ ایہ صرف
ایہ آکھدا اے کہ اک «کوڈ کھولّن» والا فنکشن موجود اے جیہدے لئی
ہر سلسلے دا کوئی نہ کوئی «کوڈ» موجود اے۔

$\beta$ دی لکھت گوڈل توں آئی اے۔ فیر دھیان دیو کہ معاون قضیہ
ثابت کرن دا اوکھا کم بظاہر گھٹ وسائلاں نال مناسب~$\beta$
دی تعریف کرنا اے۔ ساڈے کول
صرف ترکیب تے کمترین قدر دی تلاش دے طریقے نیں، تے بنیادی فنکشناں
وچ جمع، ضرب تے~$\Char{=}$ وی ورت سکدے آں۔ ایس معاون قضیے دے کئی
ثبوت نیں۔ اج وی سبھ توں صاف طریقیاں وچوں اک گوڈل دا اپنا طریقہ
اے۔ اوس نے نظریۂ اعداد دا اک نتیجہ ورتیا جیہنوں سن زی دا قضیہ
آکھدے نیں (روایتی ناں «چینی باقیات دا قضیہ» اے)۔

\begin{defn}
دو قدرتی عدد $a$ تے $b$ اُدوں تے صرف اُدوں \emph{آپس وچ اوّل}
ہوندے نیں جدوں اوہناں دونوں نوں پورا ونڈن والا سبھ توں وڈا
عدد~$1$ ہووے۔ یعنی اوہناں دونوں نوں پورا ونڈن والا ہور کوئی
عدد نہیں۔
\end{defn}

\begin{defn}
قدرتی عدد $a$ تے $b$ \emph{مودولو~$c$ ہم باقی} ہوندے نیں،
$a \equiv b \mod c$، اُدوں تے صرف اُدوں جدوں $c \mid (a-b)$۔
یعنی $a$ تے $b$ نوں~$c$ نال ونڈیے تاں دونوں وار اکّو باقی
بچدی اے۔
\end{defn}

سن زی دا قضیہ ایہ اے:
\begin{thm}
فرض کرو $x_0$، \dots،~$x_n$ وچ ہر جوڑے دے عدد آپس وچ اوّل
نیں، تے $y_0$، \dots،~$y_n$ کوئی وی عدد نیں۔ تاں کوئی
عدد $z$ ایسا موجود اے کہ
\begin{align*}
z & \equiv y_0 \mod x_0 \\
z & \equiv y_1 \mod x_1 \\
& \vdots  \\
z & \equiv y_n \mod x_n.
\end{align*}
\end{thm}

سن زی دا قضیہ اسیں ایویں ورتاں گے۔ جے $x_0$، \dots،~$x_n$
بالترتیب $y_0$، \dots،~$y_n$ توں وڈے ہون، تاں $z$ نوں
سلسلے $\tuple{y_0, \dots,y_n}$ دا کوڈ بنا سکدے آں۔~$y_i$
کڈھن لئی $z$ نوں~$x_i$ نال ونڈ کے باقی لینی اے۔ ایس طریقے
لئی سانوں $x_0$، \dots، ~$x_n$ دیاں مناسب قدراں لبھنیاں پین گیاں۔

ایس کم وچ دو گلاں ساڈی مدد کرن گیاں۔
$y_0$، \dots،~$y_n$ دتے ہون،
تاں ایہ من لو:
\begin{align*}
j &= \max(n, y_0 + 1, \dots, y_n + 1), \\
m &= \lcm(1,\dots,j),
\end{align*}
تے ایہ من لو:
\begin{align*}
x_0 & = 1 + m \\
x_1 & = 1 + 2 \cdot m \\
x_2 & = 1 + 3 \cdot m \\
& \vdots  \\
x_n & = 1 + (n+1) \cdot m
\end{align*}
تاں دو گلّاں سچ نیں:
\begin{enumerate}
\item\ollabel{rel-prime} $x_0,\dots,x_n$ وچ ہر جوڑے دے عدد آپس وچ اوّل نیں۔
\item\ollabel{less} ہر $i$ لئی $y_i < x_i$۔
\end{enumerate}
\olref{rel-prime} لئی دو وکھ اشاریے لو۔ جے $p$ اولی عدد اے
تے $p \mid x_i$ تے $p \mid x_k$، تاں
$p \mid 1 + (i+1) m$ تے
$p \mid 1 + (k+1) m$۔
پر فیر $p$ اوہناں دے فرق
نوں پورا ونڈدا اے:
\[
(1 + (i+1)m) - (1+ (k+1)m) = (i-k) m.
\]
کیوں جے $p$، $1 + (i+1)m$ نوں پورا ونڈدا اے، ایس لئی
$m$ نوں پورا نہیں ونڈ سکدا۔ جے ونڈدا ہووے تاں پہلی ونڈ وچ
باقی~$1$ بچدی۔ ایس لئی $p$، $i-k$ نوں پورا ونڈدا اے، کیوں
جے $p$، $(i-k)m$ نوں پورا ونڈدا اے۔ پر $\left|i-k\right|$
ودھ توں ودھ~$n$ اے، تے اسیں $j \geq n$ چُنیا اے۔ ایس توں
$p \mid m$ نکلدا اے، جیہڑا تضاد اے۔ سو کوئی اولی عدد $x_i$
تے $x_k$ دونوں نوں پورا نہیں ونڈدا۔ شرط~\olref{less} سِدھی اے:
$y_i < j \leq m < x_i$۔

ہن $\beta$ فنکشن دا معاون قضیہ ثابت کریے۔ یاد رکھو کہ اسیں $0$،
جانشین، جمع، ضرب، $\Char{=}$ تے پروجیکشن ورت سکدے آں۔ اوہناں
توں ترکیب تے باقاعدہ فنکشناں اُتے کمترین قدر دی تلاش نال بنیا
کوئی وی فنکشن ورت سکدے آں۔ جے کسے تعلق دے امتیازی فنکشن دی
تعریف ایویں ہو سکے تاں اوہ تعلق وی ورت سکدے آں۔ پہلے وانگ
وکھا سکدے آں کہ بولیائی عملاں تے محدود مسوّریت نال بنے تعلق
وی ایس جماعت وچ رہندے نیں۔ مثال لئی:
\begin{align*}
\fn{not}(x) & \defis \Char{=}(x,0)\\
\bmin{x \leq z}{R(x,y)} & \defis \umin{x}{(R(x,y) \lor x = z)}\\
\bexists{x \leq z}{R(x,y)} & \defiff R(\bmin{x \leq z}{R(x,y)}, y)
\end{align*}
فیر اسیں دکھا سکدے
آں کہ ہِٹھلیاں
ساریاں شےواں وی
بنیادی بازگشت توں
بغیر تعریف ہو
سکدیاں نیں:
\begin{enumerate}
\item جوڑی فنکشن،
$J(x,y) = \frac{1}{2}[(x+y)(x+y+1)] + x$؛
% شاید ایتھے ہور دَسّنا چاہیدا اے کہ کیہ ہو رہیا اے؛گل کجھ الجھی اے۔
\item پروجیکشن فنکشن
\begin{align*}
K(z) & = \bmin{x \leq z}{\bexists{y \leq z}{z = J(x,y)}},\\
L(z) & = \bmin{y \leq z}{\bexists{x \leq z}{z = J(x,y)}};
\end{align*}
\item چھوٹا ہون دا تعلق $x < y$؛
\item پورا ونڈن دا تعلق $x \mid y$؛
% \item $x \tsub y$
% \item $\fn{Prime}(x)$
% \item فرض کرو $p$ اولی اے؛تعلق «$x$،$p$ دی اک قوت اے»:
% \[
% \bforall{y \leq x}{(y \mid x \lif y = 1 \lor y = x)}.
% \]
\item فنکشن $\fn{rem}(x,y)$، جیہڑا $y$ نوں~$x$ نال ونڈن توں بچی
  باقی دیندا اے۔
\end{enumerate}
ہن ایہ تعریف کرو:
\begin{align*}
\beta^*(d_0,d_1,i) & = \fn{rem}(1+(i+1) d_1,d_0) \text{ تے}\\
\beta(d,i) & = \beta^*(K(d),L(d),i).
\end{align*}
ایہو فنکشن سانوں چاہیدا اے۔
$a_0,\dots,a_n$ اُتّے وانگ
دتے ہون،تاں ایہ من لو:
\[
j = \max(n,a_0+1,\dots,a_n+1),
\]
تے $d_1 = \lcm(1,\dots,j)$
من لو۔ اُتّے دی
\olref{rel-prime}
مطابق سانوں پتا
اے کہ $1+d_1$،
$1+2 d_1$، \dots،
$1+(n+1) d_1$ وچ ہر جوڑے دے عدد آپس وچ اوّل نیں۔ نالے
~\olref{less} مطابق اوہ سارے $a_0,\dots,a_n$ توں وڈے نیں۔
سن زی دے قضیے
مطابق کوئی قدر~$d_0$ موجود اے جیہڑی ہر~$i$ لئی ایہ شرط پوری کردی اے۔
\[
d_0 \equiv a_i \mod (1+(i+1)d_1)
\]
تے سو (کیوں جے
$d_1$،~$a_i$
توں وڈّا اے):
\[
a_i = \fn{rem}(1+(i+1)d_1,d_0).
\]
$d = J(d_0,d_1)$
من لو۔ تاں ہر
$i \le n$
لئی ساڈے کول
ایہ اے:
\begin{align*}
\beta(d,i) & =  \beta^*(d_0,d_1,i) \\
& =  \fn{rem}(1+(i+1) d_1,d_0) \\
& =  a_i
\end{align*}
ایہو سانوں چاہیدا
سی۔ ایس نال
$\beta$-فنکشن دے
معاون قضیے دا
ثبوت مکمل اے۔

\begin{prob}
  وکھاؤ کہ تعلق $x < y$، $x \mid y$ تے فنکشن~$\fn{rem}(x,y)$
دیاں تعریفاں بنیادی بازگشت توں بغیر دِتیاں جا سکدیاں نیں۔ تسیں
$0$، جانشین، جمع، ضرب، $\Char{=}$، پروجیکشن، محدود کمترین قدر
دی تلاش تے محدود مسوّریت ورت سکدے او۔
\end{prob}

\end{document}
